\begin{tikzpicture}
  \foreach \i in {0,...,3}{
    \node[flipflop T] (ff\i) at (\i*4.2,0) {};
    \node[below=1mm of ff\i.south, font=\small] {FF$_{\i}$};
    \draw (ff\i.pin 6) -- ++(0.4,0) coordinate (q\i) -- ++(0,2.6) node[above] {$Q_{\i}$};
    \fill (q\i) circle (1.3pt);
    \draw (ff\i.pin 3) -- ++(-0.35,0) -- ++(0,-1.4) coordinate (k\i);
  }
  \draw (ff0.pin 1) -- ++(-0.6,0) node[left] {1};
  \draw (q0 |- ff1.pin 1) -- (ff1.pin 1);
  \fill (q0 |- ff1.pin 1) circle (1.3pt);
  \node[and port, scale=0.8] (a2) at ($(ff2.pin 1)+(-1.1,1.3)$) {};
  \node[and port, scale=0.8] (a3) at ($(ff3.pin 1)+(-1.1,1.3)$) {};
  \draw (a2.out) -- ++(0.2,0) |- (ff2.pin 1);
  \draw (a3.out) -- ++(0.2,0) |- (ff3.pin 1);
  \draw (q0 |- a2.in 1) -- (a2.in 1);  \fill (q0 |- a2.in 1) circle (1.3pt);
  \draw (q1 |- a2.in 2) -- (a2.in 2);  \fill (q1 |- a2.in 2) circle (1.3pt);
  \draw ($(a2.out)+(0.2,0)$) |- ($(a3.in 1)+(0,0.0)$) ;
  \fill ($(a2.out)+(0.2,0)$) circle (1.3pt);
  \draw (q2 |- a3.in 2) -- (a3.in 2);  \fill (q2 |- a3.in 2) circle (1.3pt);
  \draw[acento, line width=1pt] (k0) -- (k3);
  \draw[acento, line width=1pt] (k0) -- ++(-0.8,0) node[left, text=tinta] {clk};
  \foreach \i in {0,...,3}{\fill[acento] (k\i) circle (1.3pt);}
  \node[nota, anchor=north west, text width=15cm] at ($(k0)+(-1.4,-0.3)$) {$T_0 = 1$, $T_1 = Q_0$, $T_2 = Q_1Q_0$, $T_3 = Q_2Q_1Q_0$. Todos los flip-flops reciben el mismo reloj (en azul): por eso es sincrónico.};
\end{tikzpicture}
